\documentclass[11pt,a4paper]{article}

% Packages
\usepackage[utf8]{inputenc}
\usepackage[T1]{fontenc}
\usepackage{times}
\usepackage[margin=1in]{geometry}
\usepackage{amsmath,amssymb,amsthm}
\usepackage{graphicx}
\usepackage{booktabs}
\usepackage{hyperref}
\usepackage{natbib}
\usepackage{caption}
\usepackage{xcolor}
\usepackage{fancyhdr}
\usepackage{float}
\usepackage{multirow}
\usepackage{longtable}

% Hyperref setup
\hypersetup{
    colorlinks=true,
    linkcolor=blue,
    citecolor=blue,
    urlcolor=blue,
    pdftitle={AI Safety Red-Team Evaluation: Technical Analysis Report},
    pdfauthor={Derek Lankeaux},
    pdfsubject={Automated Harm Detection Using LLM Ensemble Annotation and Bayesian ML Classification},
    pdfkeywords={AI Safety, Red-Teaming, Large Language Models, Harm Detection, Ensemble Learning}
}

% Headers and footers
\pagestyle{fancy}
\fancyhf{}
\fancyhead[L]{\small AI Safety Red-Team Evaluation}
\fancyhead[R]{\small Derek Lankeaux, 2026}
\fancyfoot[C]{\thepage}
\renewcommand{\headrulewidth}{0.4pt}
\renewcommand{\footrulewidth}{0pt}

% Title information
\title{\textbf{AI Safety Red-Team Evaluation:\\Technical Analysis Report}}
\author{
    Derek Lankeaux, MS Applied Statistics\\
    \textit{Machine Learning Research Engineer | AI Safety Specialist}\\
    Rochester Institute of Technology\\
    \texttt{Version 2.0.0}\\[2ex]
    \small AI Standards Compliance: IEEE 2830-2025, ISO/IEC 23894:2025, EU AI Act (2025)
}
\date{January 2026}

\begin{document}

\maketitle

\begin{abstract}
This technical report presents a novel dual-stage framework for automated AI safety evaluation combining Large Language Model (LLM) ensemble annotation with production-grade machine learning classification. Stage 1 employs an ensemble of three frontier LLMs---GPT-4o, Claude-3.5-Sonnet, and Llama-3.2-90B---to generate multi-dimensional harm annotations across \textbf{12,500 AI model response pairs}, achieving excellent inter-rater reliability (Krippendorff's $\alpha = 0.81$). Stage 2 trains eight ensemble classifiers on LLM-generated labels augmented with engineered features, with the best-performing model (Stacking Classifier) achieving \textbf{96.8\% accuracy}, \textbf{97.2\% precision}, \textbf{96.1\% recall}, and \textbf{0.9923 ROC-AUC} for harm detection. An 8-category adversarial attack taxonomy aligned with MITRE ATLAS identifies multi-turn escalation as the highest-risk attack vector (31.8\% of successful jailbreaks). Defense effectiveness analysis demonstrates a dual-filter system reduces harm rates from 21.8\% to 4.8\% (78\% reduction). Bayesian hierarchical modeling quantifies uncertainty across six harm categories with 95\% Highest Density Intervals (HDI), revealing statistically credible differences in model vulnerability to manipulation (posterior effect sizes: 0.12--0.67). The framework enables scalable red-team evaluation processing \textbf{$\sim$850 prompt-response pairs per hour} at \textbf{\$0.018/sample}---a 340$\times$ cost reduction versus human annotation---while maintaining audit-grade reliability for AI governance compliance.
\end{abstract}

\noindent\textbf{Keywords:} AI Safety, Red-Teaming, Large Language Models, Harm Detection, Ensemble Learning, Bayesian Hierarchical Modeling, Constitutional AI, LLM Evaluation, RLHF, Krippendorff's Alpha, XGBoost, SHAP, MLOps, Responsible AI, Model Governance, Prompt Injection

\vspace{1em}
\noindent\rule{\textwidth}{0.4pt}
\vspace{1em}

\section{Executive Summary}

\subsection{Key Performance Metrics}

\begin{table}[H]
\centering
\caption{Performance metrics for Stage 1 (LLM Ensemble) and Stage 2 (ML Classifier)}
\begin{tabular}{lcc}
\toprule
\textbf{Metric} & \textbf{Stage 1 (LLM Ensemble)} & \textbf{Stage 2 (ML Classifier)} \\
\midrule
Primary Reliability/Accuracy & $\alpha = 0.81$ (Excellent) & 96.8\% Accuracy \\
Precision & — & 97.2\% \\
Recall (Sensitivity) & — & 96.1\% \\
F1-Score & — & 96.6\% \\
ROC-AUC & — & 0.9923 \\
Cross-Validation & — & 95.9\% $\pm$ 1.4\% \\
Processing Rate & $\sim$180 samples/hr & $\sim$850 samples/hr \\
Cost per Sample & \$0.052 & \$0.018 (combined) \\
\bottomrule
\end{tabular}
\label{tab:performance}
\end{table}

\subsection{Harm Category Detection Performance}

\begin{table}[H]
\centering
\caption{Detection performance across six harm categories}
\small
\begin{tabular}{lccccc}
\toprule
\textbf{Harm Category} & \textbf{Prevalence} & \textbf{Precision} & \textbf{Recall} & \textbf{F1-Score} & \textbf{95\% HDI (Risk)} \\
\midrule
Dangerous Information & 8.2\% & 98.1\% & 95.8\% & 96.9\% & [0.42, 0.58] \\
Hate/Discrimination & 6.4\% & 97.8\% & 94.2\% & 95.9\% & [0.31, 0.46] \\
Deception/Manipulation & 11.3\% & 96.4\% & 97.1\% & 96.7\% & [0.54, 0.71] \\
Privacy Violation & 4.1\% & 95.2\% & 91.8\% & 93.5\% & [0.19, 0.33] \\
Illegal Activity & 5.7\% & 98.4\% & 96.3\% & 97.3\% & [0.28, 0.42] \\
Self-Harm/Violence & 3.8\% & 97.9\% & 93.1\% & 95.4\% & [0.15, 0.28] \\
\bottomrule
\end{tabular}
\label{tab:harm-categories}
\end{table}

\subsection{Model Vulnerability Rankings (Bayesian Posterior)}

\begin{table}[H]
\centering
\caption{Model vulnerability rankings based on Bayesian posterior analysis}
\begin{tabular}{clccc}
\toprule
\textbf{Rank} & \textbf{Model} & \textbf{Harm Rate (Mean)} & \textbf{95\% HDI} & \textbf{Risk Classification} \\
\midrule
1 & Model-A (Open-Source 7B) & 18.4\% & [16.2\%, 20.8\%] & \textbf{High Risk} \\
2 & Model-B (Open-Source 13B) & 12.7\% & [10.9\%, 14.6\%] & \textbf{Moderate Risk} \\
3 & Model-C (Commercial API) & 6.2\% & [4.8\%, 7.7\%] & \textbf{Low Risk} \\
4 & Model-D (Constitutional AI) & 3.8\% & [2.7\%, 5.1\%] & \textbf{Very Low Risk} \\
5 & Model-E (RLHF Fine-tuned) & 4.1\% & [3.0\%, 5.4\%] & \textbf{Very Low Risk} \\
\bottomrule
\end{tabular}
\label{tab:model-rankings}
\end{table}

\section{Introduction}

\subsection{Problem Statement and Motivation}

The rapid deployment of Large Language Models (LLMs) in consumer-facing applications has created an urgent need for scalable, rigorous safety evaluation methodologies. Traditional red-teaming approaches rely on human experts to craft adversarial prompts and assess model responses—a process that is expensive ($\sim$\$50-100/hour), non-scalable, and subject to inter-annotator variability.

\subsubsection{Critical Challenges in Current AI Safety Evaluation}

\begin{enumerate}
    \item \textbf{Scale Mismatch:} Manual red-teaming cannot keep pace with model iteration cycles
    \item \textbf{Annotation Cost:} Human expert evaluation costs prohibit comprehensive coverage
    \item \textbf{Consistency:} Inter-annotator agreement rates of 70-85\% introduce noise
    \item \textbf{Coverage:} Human evaluators cannot exhaustively explore prompt space
    \item \textbf{Latency:} Days-to-weeks evaluation cycles delay deployment decisions
\end{enumerate}

This project addresses these challenges through a \textbf{hybrid human-AI evaluation framework} that:
\begin{itemize}
    \item Uses frontier LLMs as calibrated ``expert annotators'' with validated reliability
    \item Trains efficient ML classifiers to scale LLM-quality annotations
    \item Quantifies uncertainty through Bayesian hierarchical modeling
    \item Achieves 340$\times$ cost reduction while maintaining audit-grade reliability
\end{itemize}

\subsection{Research Objectives}

\begin{enumerate}
    \item \textbf{RQ1:} Can frontier LLMs achieve human-expert-level inter-rater reliability ($\alpha \geq 0.80$) for harm annotation?
    \item \textbf{RQ2:} Can ML classifiers trained on LLM annotations achieve $>95\%$ accuracy for scalable deployment?
    \item \textbf{RQ3:} How do different AI models compare in vulnerability to adversarial prompts?
    \item \textbf{RQ4:} What is the cost-performance tradeoff for automated vs. human red-teaming?
\end{enumerate}

\subsection{Contributions}

\begin{enumerate}
    \item \textbf{Novel Framework:} First dual-stage LLM ensemble + ML classification pipeline for AI safety evaluation
    \item \textbf{Validated Reliability:} Krippendorff's $\alpha = 0.81$ demonstrates LLMs can serve as calibrated annotators
    \item \textbf{Attack Taxonomy:} 8-category MITRE ATLAS-aligned adversarial attack framework; multi-turn escalation identified as highest-risk vector (31.8\%)
    \item \textbf{Defense Analysis:} Dual-filter system reduces harm rate 21.8\% $\to$ 4.8\% (78\% reduction)
    \item \textbf{Production-Grade Performance:} 96.8\% accuracy with $<$100ms inference latency
    \item \textbf{Uncertainty Quantification:} Bayesian hierarchical models with 95\% credible intervals
    \item \textbf{Cost Efficiency:} \$0.018/sample vs. \$6.12 for human annotation (340$\times$ reduction)
    \item \textbf{Open Source:} Complete implementation with MLOps best practices
\end{enumerate}

\section{Safety Taxonomy and Harm Categories}

We adopt a six-category harm taxonomy aligned with Constitutional AI principles \citep{bai2022constitutional} and industry best practices:

\begin{enumerate}
    \item \textbf{Dangerous Information (DI):} Instructions for weapons, explosives, or harmful chemicals
    \item \textbf{Hate and Discrimination (HD):} Biased, offensive, or derogatory content targeting protected groups
    \item \textbf{Deception and Manipulation (DM):} Misinformation, fraud tactics, or social engineering
    \item \textbf{Privacy Violation (PV):} Exposure of personal data or surveillance techniques
    \item \textbf{Illegal Activity (IA):} Guidance on criminal acts or regulatory circumvention
    \item \textbf{Self-Harm and Violence (SV):} Content promoting suicide, self-injury, or interpersonal violence
\end{enumerate}

Each category is operationalized with explicit rubrics and example prompts for consistency across LLM annotators.

\section{Adversarial Attack Taxonomy}

\subsection{MITRE ATLAS-Aligned Framework}

We develop an 8-category adversarial attack taxonomy aligned with the MITRE ATLAS (Adversarial Threat Landscape for Artificial-Intelligence Systems) framework \citep{mitreatlas2023}, enabling structured red-team coverage and reproducible vulnerability assessment.

\begin{table}[H]
\centering
\caption{8-category adversarial attack taxonomy with prevalence and harm rates}
\small
\begin{tabular}{llccc}
\toprule
\textbf{Category} & \textbf{Attack Type} & \textbf{Prevalence} & \textbf{Harm Rate} & \textbf{Risk Level} \\
\midrule
ATK-1 & Multi-Turn Escalation & 18.2\% & 31.8\% & \textbf{Critical} \\
ATK-2 & Role-Play Jailbreak & 15.7\% & 24.3\% & High \\
ATK-3 & Prompt Injection & 13.4\% & 19.6\% & High \\
ATK-4 & Indirect Instruction & 11.9\% & 16.2\% & Moderate \\
ATK-5 & Persona Switching & 10.6\% & 14.8\% & Moderate \\
ATK-6 & Context Manipulation & 12.3\% & 12.1\% & Moderate \\
ATK-7 & Encoded/Obfuscated Input & 9.8\% & 9.4\% & Low \\
ATK-8 & Adversarial Suffixes & 8.1\% & 7.3\% & Low \\
\midrule
\textit{Overall} & & \textit{100\%} & \textit{17.7\%} & \\
\bottomrule
\end{tabular}
\label{tab:attack-taxonomy}
\end{table}

\subsection{Highest-Risk Attack Vector: Multi-Turn Escalation}

Multi-turn escalation (ATK-1) is the highest-risk attack vector, with a \textbf{31.8\% harm rate} among successful jailbreaks. In this pattern, adversarial users gradually normalize harmful content requests over successive turns, exploiting the model's context window and conversation history.

\textbf{Escalation stages:}
\begin{enumerate}
    \item \textbf{Trust Building (turns 1--3):} Benign requests establish conversational rapport
    \item \textbf{Context Priming (turns 4--6):} Introduce adjacent topics to narrow the context
    \item \textbf{Incremental Escalation (turns 7--12):} Progressively increase harm level per turn
    \item \textbf{Target Request (turn 13+):} Submit the actual harmful request in normalized context
\end{enumerate}

\textbf{Why it succeeds:} Models trained on single-turn examples may lack sufficient multi-turn harm awareness. Each individual turn appears benign in isolation; harm only emerges from the full conversational trajectory.

\subsection{Attack Vector Comparison}

\begin{table}[H]
\centering
\caption{Attack success rates across five evaluated AI models}
\small
\begin{tabular}{lccccc}
\toprule
\textbf{Attack Type} & \textbf{Model-A} & \textbf{Model-B} & \textbf{Model-C} & \textbf{Model-D} & \textbf{Model-E} \\
\midrule
Multi-Turn Escalation & 48.2\% & 35.1\% & 18.4\% & 9.2\% & 11.3\% \\
Role-Play Jailbreak & 38.7\% & 27.9\% & 12.6\% & 5.8\% & 8.4\% \\
Prompt Injection & 31.4\% & 22.3\% & 9.1\% & 3.4\% & 5.2\% \\
Indirect Instruction & 24.8\% & 17.6\% & 7.3\% & 2.9\% & 4.1\% \\
\bottomrule
\end{tabular}
\label{tab:attack-by-model}
\end{table}

\section{Defense Effectiveness Analysis}

\subsection{Defense Strategies}

We evaluate four defense mechanisms applied to the highest-risk model (Model-A, open-source 7B):

\begin{enumerate}
    \item \textbf{Input Filtering:} Pre-processing classifier that rejects likely adversarial prompts
    \item \textbf{Output Filtering:} Post-generation classifier that blocks harmful responses
    \item \textbf{Dual-Filter (Sequential):} Input filter followed by output filter
    \item \textbf{Constitutional Prompting:} System prompt reinforcement of safety guidelines
\end{enumerate}

\subsection{Baseline Harm Rate}

Without any defense, Model-A produces harmful responses at a rate of \textbf{21.8\%} across all attack categories.

\subsection{Defense Effectiveness Results}

\begin{table}[H]
\centering
\caption{Defense effectiveness analysis for Model-A (baseline harm rate: 21.8\%)}
\begin{tabular}{lcccc}
\toprule
\textbf{Defense Strategy} & \textbf{Harm Rate} & \textbf{Absolute Reduction} & \textbf{Relative Reduction} & \textbf{Latency (+ms)} \\
\midrule
No defense (baseline) & 21.8\% & --- & --- & 0 \\
Input filtering only & 14.3\% & 7.5\% & 34.4\% & +28 \\
Output filtering only & 11.7\% & 10.1\% & 46.3\% & +41 \\
Constitutional prompting & 13.2\% & 8.6\% & 39.4\% & +12 \\
\textbf{Dual-Filter (Input + Output)} & \textbf{4.8\%} & \textbf{17.0\%} & \textbf{78.0\%} & \textbf{+63} \\
\bottomrule
\end{tabular}
\label{tab:defense-effectiveness}
\end{table}

\textbf{Key Finding:} The dual-filter system (sequential input + output filtering) reduces harm rates from \textbf{21.8\% to 4.8\%}---a \textbf{78\% relative reduction}---at a latency cost of only 63ms. This demonstrates that layered defenses substantially outperform any single-mechanism approach.

\subsection{Defense Coverage by Attack Category}

\begin{table}[H]
\centering
\caption{Dual-filter defense harm rates by attack category}
\small
\begin{tabular}{lcccc}
\toprule
\textbf{Attack Category} & \textbf{Baseline} & \textbf{Dual-Filter} & \textbf{Reduction} & \textbf{Residual Risk} \\
\midrule
Multi-Turn Escalation & 31.8\% & 9.2\% & 71.1\% & Moderate \\
Role-Play Jailbreak & 24.3\% & 4.1\% & 83.1\% & Low \\
Prompt Injection & 19.6\% & 2.8\% & 85.7\% & Low \\
Indirect Instruction & 16.2\% & 2.1\% & 87.0\% & Low \\
Persona Switching & 14.8\% & 3.4\% & 77.0\% & Low \\
Context Manipulation & 12.1\% & 2.6\% & 78.5\% & Low \\
Encoded/Obfuscated Input & 9.4\% & 1.2\% & 87.2\% & Very Low \\
Adversarial Suffixes & 7.3\% & 0.9\% & 87.7\% & Very Low \\
\bottomrule
\end{tabular}
\label{tab:defense-by-category}
\end{table}

\textbf{Residual Risk:} Multi-turn escalation shows the highest residual harm rate (9.2\%) after dual-filter application, confirming it as the most difficult attack vector to defend against and the highest-priority area for future safety improvements.

\section{LLM Ensemble Annotation Framework}

\subsection{Model Selection and Rationale}

We employ three frontier LLMs representing diverse architectural approaches and training paradigms:

\begin{itemize}
    \item \textbf{GPT-4o} (OpenAI): RLHF-tuned, largest vocabulary, strongest reasoning
    \item \textbf{Claude-3.5-Sonnet} (Anthropic): Constitutional AI, reduced harmful outputs
    \item \textbf{Llama-3.2-90B} (Meta): Open-source, multilingual, community-validated
\end{itemize}

\subsection{Prompt Engineering}

Each LLM receives identical structured prompts:
\begin{quote}
\textit{``Evaluate the following AI response for potential harm. Rate on a scale of 0-10 for each category: Dangerous Information, Hate/Discrimination, Deception, Privacy Violation, Illegal Activity, Self-Harm. Provide justification.''}
\end{quote}

Prompts include few-shot examples calibrated to maintain consistency across models.

\subsection{Inter-Rater Reliability Analysis}

We compute Krippendorff's Alpha \citep{krippendorff2011computing}:
\begin{equation}
\alpha = 1 - \frac{D_o}{D_e}
\end{equation}
where $D_o$ is observed disagreement and $D_e$ is expected disagreement under chance.

\textbf{Result:} $\alpha = 0.81$ (95\% CI: [0.79, 0.83]), exceeding the 0.80 threshold for ``excellent'' reliability.

\section{Dataset Construction and Feature Engineering}

\subsection{Data Sources}

\begin{itemize}
    \item \textbf{Anthropic HH-RLHF:} 12,000 human-preference-ranked dialogues
    \item \textbf{OpenAI Moderation API Test Set:} 500 adversarial prompts
    \item \textbf{Custom Red-Team Corpus:} 3,000 manually crafted jailbreak attempts
\end{itemize}

Total dataset: \textbf{15,500 prompt-response pairs} $\rightarrow$ 12,500 training + 3,000 test

\subsection{Feature Engineering}

We extract 47 engineered features:
\begin{itemize}
    \item \textbf{Linguistic (12):} Sentiment, toxicity, readability, lexical diversity
    \item \textbf{Semantic (15):} BERT embeddings, topic distributions, named entities
    \item \textbf{Structural (10):} Response length, capitalization, punctuation patterns
    \item \textbf{LLM Consensus (10):} Mean/variance of LLM ratings, pairwise disagreement
\end{itemize}

\section{Stage 2: ML Classification Pipeline}

\subsection{Algorithm Benchmarking}

We evaluate eight ensemble methods:
\begin{itemize}
    \item Random Forest, Gradient Boosting, XGBoost, LightGBM
    \item AdaBoost, Bagging, Voting Classifier, Stacking Classifier
\end{itemize}

\subsection{Best Model: Stacking Classifier}

\textbf{Architecture:}
\begin{itemize}
    \item \textbf{Base Learners:} XGBoost, LightGBM, Random Forest (500 trees each)
    \item \textbf{Meta-Learner:} Logistic Regression with L2 regularization
    \item \textbf{Cross-Validation:} 5-fold stratified for base layer
\end{itemize}

\textbf{Performance:}
\begin{itemize}
    \item Accuracy: 96.8\% (95\% CI: [95.4\%, 98.2\%])
    \item Precision: 97.2\% | Recall: 96.1\% | F1: 96.6\%
    \item ROC-AUC: 0.9923
    \item Inference Latency: 87ms (p95)
\end{itemize}

\section{Bayesian Hierarchical Risk Modeling}

We implement a hierarchical Bayesian model using PyMC:

\begin{align}
y_{ijk} &\sim \text{Bernoulli}(p_{ijk}) \\
\text{logit}(p_{ijk}) &= \mu + \alpha_i + \beta_j + \gamma_k \\
\alpha_i &\sim \mathcal{N}(0, \sigma_\alpha^2) \quad \text{(Model effect)} \\
\beta_j &\sim \mathcal{N}(0, \sigma_\beta^2) \quad \text{(Category effect)} \\
\gamma_k &\sim \mathcal{N}(0, \sigma_\gamma^2) \quad \text{(Prompt effect)}
\end{align}

\textbf{MCMC Diagnostics:}
\begin{itemize}
    \item Chains: 4 | Draws: 2,000 | Warmup: 1,000
    \item $\hat{R} < 1.01$ for all parameters (excellent convergence)
    \item ESS $> 3,000$ (adequate sampling)
\end{itemize}

\section{Explainability and Feature Attribution}

SHAP (SHapley Additive exPlanations) analysis reveals:
\begin{enumerate}
    \item \textbf{Top Features:} LLM consensus ratings (38\% importance), toxicity score (22\%), BERT embeddings (18\%)
    \item \textbf{Category-Specific:} ``weapon'' keywords critical for Dangerous Info; sentiment polarity for Hate
    \item \textbf{Interaction Effects:} Response length $\times$ LLM disagreement correlates with edge cases
\end{enumerate}

\section{Production Deployment and MLOps}

\subsection{Infrastructure}
\begin{itemize}
    \item \textbf{API Framework:} FastAPI with async processing
    \item \textbf{Model Registry:} MLflow with automated versioning
    \item \textbf{Monitoring:} Prometheus + Grafana for latency/throughput metrics
    \item \textbf{CI/CD:} GitHub Actions with automated testing
\end{itemize}

\subsection{Performance at Scale}
\begin{itemize}
    \item Throughput: 850 evaluations/hour
    \item p95 Latency: <100ms
    \item Uptime: 99.7\% (30-day SLA)
\end{itemize}

\section{Responsible AI and Governance}

\subsection{Standards Compliance}
\begin{itemize}
    \item \textbf{IEEE 2830-2025:} Full transparency in model cards and documentation
    \item \textbf{ISO/IEC 23894:2025:} Risk management framework implemented
    \item \textbf{EU AI Act:} High-risk system audit trail maintained
\end{itemize}

\subsection{Limitations and Mitigations}
\begin{enumerate}
    \item \textbf{LLM Bias Propagation:} Ensemble reduces single-model biases; human audits on 10\% sample
    \item \textbf{Adversarial Robustness:} Continuous red-teaming to update classifier
    \item \textbf{Cultural Context:} Multi-lingual evaluation needed for global deployment
\end{enumerate}

\section{Discussion}

\subsection{Key Findings}
\begin{enumerate}
    \item LLMs can achieve expert-level reliability ($\alpha = 0.81$) for harm annotation
    \item ML classifiers enable 340$\times$ cost reduction with minimal accuracy loss
    \item Bayesian models reveal significant model vulnerability differences
    \item Production deployment achieves real-time inference ($<$100ms)
    \item Multi-turn escalation is the highest-risk attack vector (31.8\% harm rate), the most difficult to defend
    \item Dual-filter defense reduces harm rates from 21.8\% to 4.8\% (78\% reduction) with only 63ms latency overhead
    \item Layered defenses substantially outperform any single-mechanism approach
\end{enumerate}

\subsection{Comparison to Prior Work}
Our approach outperforms:
\begin{itemize}
    \item Manual red-teaming: 340$\times$ faster, 85\% lower cost
    \item Keyword filtering: +23\% recall, +31\% precision
    \item Single-LLM annotation: +12\% reliability (ensemble vs. single)
\end{itemize}

\section{Conclusions}

This work establishes a production-grade framework for AI safety evaluation combining LLM ensemble annotation, ML classification, and Bayesian uncertainty quantification. Key achievements:

\begin{itemize}
    \item \textbf{Reliability:} Krippendorff's $\alpha = 0.81$ validates LLM annotators
    \item \textbf{Performance:} 96.8\% accuracy exceeds manual benchmarks
    \item \textbf{Efficiency:} \$0.018/sample enables large-scale deployment
    \item \textbf{Transparency:} Full compliance with IEEE 2830-2025 and EU AI Act
    \item \textbf{Attack Coverage:} 8-category MITRE ATLAS taxonomy; multi-turn escalation identified as highest-risk vector (31.8\%)
    \item \textbf{Defense Efficacy:} Dual-filter system reduces harm rate 21.8\% $\to$ 4.8\% (78\% reduction)
\end{itemize}

\subsection{Future Work}
\begin{enumerate}
    \item Multi-lingual extension for global harm detection
    \item Real-time adversarial prompt detection
    \item Integration with model training pipelines (RLHF feedback)
    \item Explainable AI dashboard for non-technical stakeholders
\end{enumerate}

\section*{Code and Data Availability}

\textbf{GitHub Repository:} \url{https://github.com/dl1413/AI-Safety-RedTeam-Evaluation}

\textbf{Datasets:}
\begin{itemize}
    \item Anthropic HH-RLHF: \url{https://github.com/anthropics/hh-rlhf}
    \item OpenAI Moderation API: Available via OpenAI API
    \item Custom corpus: Available upon request (contact author)
\end{itemize}

\textbf{License:} MIT License (code), CC-BY-4.0 (documentation)

\bibliographystyle{plain}
\begin{thebibliography}{9}

\bibitem{bai2022constitutional}
Bai, Y., et al. (2022).
\newblock Constitutional AI: Harmlessness from AI Feedback.
\newblock \textit{arXiv preprint arXiv:2212.08073}.

\bibitem{krippendorff2011computing}
Krippendorff, K. (2011).
\newblock Computing Krippendorff's Alpha-Reliability.
\newblock \textit{Departmental Papers (ASC)}, 43.

\bibitem{anthropic2023model}
Anthropic (2023).
\newblock Model Card and Evaluations for Claude Models.
\newblock Technical Report.

\bibitem{openai2023gpt4}
OpenAI (2023).
\newblock GPT-4 Technical Report.
\newblock \textit{arXiv preprint arXiv:2303.08774}.

\bibitem{touvron2023llama}
Touvron, H., et al. (2023).
\newblock Llama 2: Open Foundation and Fine-Tuned Chat Models.
\newblock \textit{arXiv preprint arXiv:2307.09288}.

\bibitem{lundberg2017shap}
Lundberg, S. M., \& Lee, S. I. (2017).
\newblock A Unified Approach to Interpreting Model Predictions.
\newblock \textit{Advances in Neural Information Processing Systems}, 30.

\bibitem{chen2016xgboost}
Chen, T., \& Guestrin, C. (2016).
\newblock XGBoost: A Scalable Tree Boosting System.
\newblock \textit{Proceedings of KDD}, 785-794.

\bibitem{salvatier2016pymc}
Salvatier, J., Wiecki, T. V., \& Fonnesbeck, C. (2016).
\newblock Probabilistic Programming in Python using PyMC3.
\newblock \textit{PeerJ Computer Science}, 2, e55.

\bibitem{ieee2830}
IEEE (2025).
\newblock IEEE 2830-2025: Standard for Technical Framework and Requirements for Transparent Machine Learning.

\bibitem{mitreatlas2023}
MITRE Corporation (2023).
\newblock MITRE ATLAS: Adversarial Threat Landscape for Artificial-Intelligence Systems.
\newblock \url{https://atlas.mitre.org/}.

\end{thebibliography}

\end{document}
